\documentclass[11pt]{article}

\begin{document}

\title{McPhase Explorer}
\author{Till Hoffmann}
\maketitle

\section{Introduction}
\emph{McPhase Explorer} is an IDE - like program based on the NetBeans platform (http://www.platform.netbeans.org/) designed to help users create, edit and view files associated with the \emph{McPhase} program package.

Since the program is based on a well established platform, documentation to add and/or change the program's behaviour is readily available at http://www.platform.netbeans.org/kb/.

\section{Architecture}
The NetBeans platform uses a modular architecture. E.g. each file type support is capsuled in its own independent module which can be loaded in any NetBeans platform application. To get a quick overview of the architecture please visit http://platform.netbeans.org/tutorials/nbm-quick-start.html.

\subsection{McPhaseBase}
\emph{McPhaseBase} is a module containing basic implementations of commonly used classes to streamline the development process of file type support modules. It provides
\begin{itemize}
\item support for multiple editors being associated with one file such as a plain text representation and a UI based representation,
\item a basic implementation of a plain text editor (\emph{SourceEditorBase}),
\item a base class for visual editors (\emph{VisualEditorBase}) including support for getting and setting key-value type variables.
\end{itemize}

\subsection{Adding file type support}
Adding editor support for a custom file type can be achieved easily by 
\begin{enumerate}
\item adding a new module to the \emph{McPhaseExplorer} application suite,
\item choosing \emph{Module Development $\rightarrow$ File Type} in the new file wizard,
\item letting the newly created \emph{DataObject} inherit from \emph{org.mcphase.DataObjectBase} and implementing its abstract methods.
\end{enumerate}
The only abstract method of \emph{org.mcphase.base.DataObjectBase} is \emph{DataEditorSupport createEditorSupport()} and returns an editor support for the file type. Custom editor supports can be provided (http://bits.netbeans.org/dev/javadoc/org-openide-loaders/org/openide/text/DataEditorSupport.html) or the helper class \emph{org.mcphase.base.EditorSupportBase} can be used as shown in the following sample implementation
\begin{verbatim}
[...]
public DataEditorSupport createEditorSupport() {
    //Create a new editor support associated with this data object
    EditorSupportBase support = new EditorSupportBase(this);
    //Create a view description
    ViewBase visual = new ViewBase();
    //Set its display name
    visual.setDisplayName("some name"));
    //Set its icon (extending Image)
    visual.setIcon(someIcon));
    //Set its editor (implementing MultiViewElement 
    //often extending VisualEditorBase)
    visual.setViewElement(someEditor);
    //Create the view elements
    MultiViewDescription[] desc = {new SourceView(support), visual};
    //Add the view elements to the editor support
    support.setViewDescriptions(desc);
    //Set the default element
    support.setDefaultViewDescription(desc[0]);
    //return the editor support
    return support;
    }
[...]
\end{verbatim}
For a full implementation see the sample implementation of mcphas.ini. All classes used above are documented in either of the following javadoc documentations
\begin{description}
\item[NetBeans API List (http://bits.netbeans.org/dev/javadoc/index.html)]
DataEditorSupport, MultiViewDescription
\item[McPhaseBase Documentation (provided with this document)]
EditorSupportBase, SourceView, ViewBase, VisualEditorBase
\end{description}

\section{Current development status}
Internal support for the following file types is provided
\begin{description}
\item[sipf]Single ion parameter file
\item[ini]McPhase algorithm control file
\item[jpg, jpeg, png, gif, bmp] Image files
\item[ant] Ant build scripts
\end{description}

All other file types are opened with the native OS application.

\section{McPhaseBase Documentation}

\subsection{org.mcphase.base.DataObjectBase}
The following is a brief summary of the most commonly used classes and their interfaces. A full documentation can be found in the javadoc provided with this document.
\emph{org.mcphase.base.DataObjectBase} represents a file object and provides support for multiple editors being associated with it. 

It has one abstract method, \emph{DataEditorSupport createEditorSupport()}, which should be overridden in derived classes to provide an editor support. The helper class \emph{org.mcphase.base.EditorSupportBase} may be helpful in doing so.

\subsection{org.mcphase.base.EditorSupportBase}
\emph{org.mcphase.base.EditorSupportBase} implements a basic editor support to be used with multiple editors.

It requires a DataObject to be passed to its constructor and provides the following functionality
\begin{description}
\item[MultiViewDescription'' getViewDescriptions()]
Gets the view descriptions representing the various editor associated with the editor support.
\item[void setViewDescriptions(MultiViewDescription'' viewDescriptions)]
Sets the view descriptions representing the various editor associated with the editor support.
\item[MultiViewDescription getDefaultViewDescription()]
Gets the view description representing the editor shown by default.
\item[void setDefaultViewDescription(MultiViewDescription defaultViewDescription)]
Sets the view description representing the editor shown by default.
\end{description}

\subsection{org.mcphase.base.SourceEditor}
\emph{org.mcphase.base.SourceEditor} implements an editor providing access to a file using a plain text interface. Please see 'Rich Client Programming: Plugging into the NetBeans Platform' for a full description of the implementation.

\subsection{org.mcphase.base.SourceView}
\emph{org.mcphase.base.SourceView} implements a view description for a plain text editor.

\subsection{org.mcphase.base.ViewBase}
\emph{org.mcphase.base.ViewBase} implements a MultiViewDescription to streamline the development process.

It provides the following functionality
\begin{description}
\item[MultiViewElement getViewElement()]
Gets the element associated with the view (i.e. the editor).
\item[void setViewElement(MultiViewElement viewElement)]
Sets the element associated with the view (i.e. the editor).
\item[String getDisplayName()]
Gets the name displayed in the tab bar.
\item[void setDisplayName(String displayName)]
Sets the name displayed in the tab bar.
\item[Image getIcon()]
Gets the icon displayed next to the file name while the view is selected.
\item[void setIcon(Image icon)]
Sets the icon displayed next to the file name while the view is selected.
\end{description}

\subsection{org.mcphase.base.VisualEditorBase}
\emph{org.mcphase.base.VisualEditorBase} implements a basic visual editor, which can be embedded in a multiple-editor environment. Components may be registered with key-value type variables such that they are kept synchronised.

It provides the following functionality
\begin{description}
\item[void updateControls()]
Updates the content of all registered controls.
\item[void registerComponent(String variable, Component component)]
Registers the specified component with the specified variable such that they are kept synchronised. Supported components are JCheckBox, JComboBox, JTextComponent.
\item[String getVariable(String name)]
Gets the value of the named variable or \emph{null} if it cannot be found.
\item[void setVariable(String variable, String value)]
Sets the value of the named variable or removes the variable if the value is \emph{null}.
\end{description}

\end{document}